% =====================================================
% 题型：余弦定理求角选择题 (q_tjw_20260606_19_cosine_rule_cos_angle.tex)
% =====================================================
% =====================================================
% 难度：★★ (中档)
% =====================================================
\newcommand{\qTriangleCosineRuleCosStem}{%
  在 \(\triangle ABC\) 中，\(\cos C = \dfrac{2}{3}\)，\(AC = 4\)，\(BC = 3\)，则 \(\cos B =\) \hfill （\quad）%
}

\newcommand{\qTriangleCosineRuleCosOptions}{%
  \mcoptions[4]
    {\(\dfrac{1}{9}\)}
    {\(\dfrac{1}{3}\)}
    {\(\dfrac{1}{2}\)}
    {\(\dfrac{2}{3}\)}%
}

\newcommand{\qTriangleCosineRuleCosAnswer}{A}

\newcommand{\qTriangleCosineRuleCosSolution}{%
  设 \(AB=c\)。由余弦定理得
  \[
    AB^2
    = AC^2 + BC^2 - 2 \cdot AC \cdot BC \cos C
    = 16 + 9 - 2 \cdot 4 \cdot 3 \cdot \frac{2}{3}
    = 9.
  \]
  所以 \(AB=3\)。

  再由余弦定理得
  \[
    \cos B
    = \frac{AB^2 + BC^2 - AC^2}{2 \cdot AB \cdot BC}
    = \frac{9 + 9 - 16}{2 \cdot 3 \cdot 3}
    = \frac{1}{9}.
  \]
  故选 A。%
}

\newcommand{\qTriangleCosineRuleCosDiagnosis}{%
  （留白待收集学生作答后补充）%
}

\newcommand{\qTriangleCosineRuleCosTeachingNotes}{%
  快讲候选。核心是“先用夹角 \(C\) 求对边 \(AB\)，再换角 \(B\) 写余弦定理”。易错点是把 \(AC,BC\) 误当作角 \(B\) 的邻边直接代入。%
}


% =====================================================
% 别名兼容：旧课次宏定义
% =====================================================
\newcommand{\qTJWZeroSixNineteenStem}{\qTriangleCosineRuleCosStem}
\newcommand{\qTJWZeroSixNineteenOptions}{\qTriangleCosineRuleCosOptions}
\newcommand{\qTJWZeroSixNineteenAnswer}{\qTriangleCosineRuleCosAnswer}
\newcommand{\qTJWZeroSixNineteenSolution}{\qTriangleCosineRuleCosSolution}
\newcommand{\qTJWZeroSixNineteenDiagnosis}{\qTriangleCosineRuleCosDiagnosis}
\newcommand{\qTJWZeroSixNineteenTeachingNotes}{\qTriangleCosineRuleCosTeachingNotes}
